\section{Ojou-sama: статус, манеры, дистанция, "ho-ho-ho", власть через воспитание}

Ojou-sama — статус как визитная карточка: дистанция, манеры, воспитание и фирменный смех "хо-хо-хо", который сигнализирует не веселье, а иерархию. Кагуя Синомия из "Kaguya-sama: Love Is War" — точный современный пример именно потому, что сериал прямо тематизирует эту механику: её статус и самообладание — не просто фон, а оружие в игре, где проиграть значит первой признаться в чувствах.

Барьер здесь родственен himedere (3.7), но не тождественен ему: himedere требует признания статуса как условия для тепла, ojou-sama как амплуа шире — это в принципе способ держать дистанцию через воспитание, необязательно завязанный на романтическую формулу. Механика похожая: дистанция не случайна, а выучена, как оннагата в 1.1, — код манер, который персонаж исполняет, а не просто "имеет".

И здесь тот же урок 4.2, только с другой стороны: Кагуя работает как персонаж не потому, что безупречно держит статус, а потому что статус у неё систематически рушится в моменты, которые сериал делает комическими, а не унизительными, — она гениальна почти во всём и катастрофически беспомощна именно в том единственном разговоре, ради которого всё затевалось. Баг не противоречит статусу — он находится ровно в его слепом пятне.

Для Рин ojou-sama-дистанция как ядро не нужна (это снова тянет к kamidere-подобному превосходству, запрещённому в 3.8), но техника "дистанция через код, а не через прямой отказ" полезна: Рин может держать границу манерой, тоном, точностью формулировки — не говоря прямо "не подходи", а делая грамматику холода видимой. И если где-то читатель увидит трещину в этой дистанции, трещина не должна быть эффектной — она должна быть немного нелепой, как у Кагуи, а не только трогательной.
